\documentclass[11pt]{article}
\usepackage[margin=25mm]{geometry}
\usepackage{fontspec}
\setmainfont{Latin Modern Roman}
\usepackage{amsmath,amssymb,mathtools}
\usepackage{unicode-math}
\setmathfont{Latin Modern Math}
\usepackage{xurl,hyperref}
\hypersetup{colorlinks=true,urlcolor=blue}
\setcounter{secnumdepth}{0}
\setlength{\parindent}{0pt}
\setlength{\parskip}{5pt}
\setlength{\emergencystretch}{3em}
\newcommand{\sref}[2]{\hyperlink{source:#1:#2}{[S#2]}}
\newcommand{\cataloguescope}[1]{\paragraph{Recorded target.}#1}
\begin{document}
\section{Brocard's Conjecture (Prime Gaps)}\label{statement}

\subsection{Definitions and mathematical statement}

\paragraph{Shared notation.}$\mathbb N=\{1,2,\ldots\}$, and $\mathbb Z,\mathbb Q,\mathbb R,\mathbb C$ denote the integers, rationals, reals and complex numbers. Any other conventions are specified locally.


Let $p_n$ be the $n$th prime in increasing order, so $p_1=2$. Write $\pi(x)=\#\{p\leq x:p\text{ is prime}\}$. The standard Brocard assertion excludes $n=1$, as the catalogue background specifies: $(p_1^2,p_2^2)=(4,9)$ contains exactly the two primes $5$ and $7$. Endpoints $p_n^2$ are not prime.
\paragraph{Statement recorded in the supplied source.}
For every integer $n\geq2$, are there at least four primes in the open interval $(p_n^2,p_{n+1}^2)$? Equivalently, is
\[
\pi(p_{n+1}^2)-\pi(p_n^2)\geq4\qquad(n\geq2)?
\] \sref{statement}{2}

\subsection{Short English statement}
Beyond the first pair of primes, do squares of consecutive primes always enclose at least four primes?
\subsection{Sources}
\begin{description}
\item[\hypertarget{source:statement:1}{S1}] ULAM.AI and the UnsolvedMath Contributors, UnsolvedMath: A Curated Collection of Open Mathematics Problems, catalogue entry. CC BY 4.0. \url{https://huggingface.co/datasets/ulamai/UnsolvedMath}.
\item[\hypertarget{source:statement:2}{S2}] Germ\'{a}n Andr\'{e}s Paz, \emph{On Legendre\textquotesingle s, Brocard\textquotesingle s, Andrica\textquotesingle s, and Oppermann\textquotesingle s Conjectures}, arXiv:1310.1323v2 (2014), Section 3. \url{https://arxiv.org/pdf/1310.1323}.
\end{description}
\label{end:statement}
\end{document}
